
\enlargethispage{.5cm}
\skpt
\begingroup
\shorthandoff{<>}% Babel español trata de otro modo los delimitadores de curvatura de Xy-pic.
\setcounter{secnumdepth}{3}
\setcounter{subsubsection}{0}
\renewcommand{\thesubsubsection}{\arabic{section}.\arabic{proposition}.\arabic{subsubsection}}
\begin{definitions}
\label{XIII.1.10}

\subsubsection{}
\label{XIII.1.10.1}
Sea $E$ una categoría y consideremos un diagrama
$$
\xymatrix@C=.6cm{ \Phi \ar[r]^p & \Phi_1 \ar@<2pt>[r]^{p_1}
\ar@<-2pt>[r]_{p_2} & \Phi_2 } \text{,}
$$
donde $\Phi$, $\Phi_1$, $\Phi_2$ son categorías fibradas sobre~$E$ y las
flechas son morfismos de categorías fibradas, y sea
$$
a \colon p_1 p \isomto p_2 p
$$
un isomorfismo de funtores.

Se dice que el diagrama anterior es exacto si se satisface la condición
siguiente:

a) Para todo par de objetos $x$, $y$ de~$\Phi$ y todo morfismo $f
\colon p(x) \to p(y)$ tal que $p_1(f) = p_2(f)$ (identificando $p_1p$ y
$p_2p$ mediante~$a$), existe un único morfismo $g \colon x \to y$ tal que
$p(g)=f$.

\subsubsection{}
\label{XIII.1.10.2}
Consideremos el diagrama
\ifthenelse{\boolean{orig}}
{$$
\xymatrix{ \Phi \ar[r]^p & \Phi_1 \ar@<0.6ex>[r]^{p_1}
\ar@<-0.6ex>[r]_{p_2} & \Phi_2 \ar@<2ex>[r]^{p_{23}} \ar[r]^{p_{31}}
\ar@<-0.6ex>[r]_{p_{12}} & \Phi_3 } \text{,}
$$}
{$$
\xymatrix@C=.5cm{ \Phi \ar[r]^p & \Phi_1 \ar@<0.6ex>[rr]^{p_1}
\ar@<-0.6ex>[rr]_{p_2} && \Phi_2 \ar@<2ex>[rrr]^{p_{23}} \ar[rrr]^{p_{31}}
\ar@<-0.6ex>[rrr]_{p_{12}} &&& \Phi_3 } \text{,}
$$}
donde
\marginpar{358}
$\Phi$, $\Phi_i$, $1 \leqslant i \leqslant 3$, son categorías fibradas
sobre~$E$ y las flechas son morfismos de categorías fibradas. Supongamos
dados unos isomorfismos de funtores
$$
a \colon p_1 p \isomto p_2 p
$$
$$
a_1 \colon p_{31} p_2 \isomto p_{12} p_1 \text{,}\quad a_2 \colon p_{12}
p_2 \isomto p_{23} p_1 \text{,}\quad a_3 \colon p_{23} p_2 \isomto
p_{31}p_1
$$
tales que el diagrama siguiente sea conmutativo:
$$
\xymatrix@C=1.2cm{ p_{23} p_1 p \ar[r]^{\id.a} & p_{23} p_2 p \ar[r]^{a_3.\id}
& p_{31} p_1 p \ar[d]^{\id.a} \\ p_{12} p_2 p \ar[u]^{a_2.\id} &
p_{12} p_1 p \ar[l]^{\id.a} & p_{31} p_2 p \ar[l]^{a_1.\id} } \text{.}
$$
Identificamos $p_1 p$ con $p_2 p$, $p_{31} p_2$ con $p_{12} p_1$, etc.

Se dice que el diagrama anterior es exacto si se satisfacen las condiciones
siguientes:
\begin{enumerate}
\item[a)] Análoga a la condición a) de~\Ref{XIII.1.10.1}.
\item[b)] Para todo objeto $x_1$ de~$\Phi_1$ y todo isomorfismo
$u \colon p_1(x_1) \isomto p_2(x_1)$ tal que
\begin{equation*}
\label{eq:XIII.1.10.2.1}
\tag{\thesubsubsection.1} { \ifthenelse{\boolean{orig}}
{p_{23}(u) p_{31}(u)}
{p_{31}(u) p_{23}(u)} = p_{12}(u)^{-1}
\text{,} }
\end{equation*}
existe un objeto $x$ de~$\Phi$ y un isomorfismo $i
\colon p(x) \isomto x_1$ que hace conmutativo el diagrama
\begin{equation*}
\label{eq:XIII.1.10.2.2}
\tag{\thesubsubsection.2}
\begin{array}{c}
{ \xymatrix{ p_1p(x) \ar@{=}[r]
\ar[d]_{p_1(i)} & p_2p(x) \ar[d]^{p_2(i)} \\ p_1(x_1) \ar[r]^{\sim}_u
& p_2(x_1) } }
\end{array}
\end{equation*}
\end{enumerate}

\subsubsection{}
\label{XIII.1.10.3}
Se define de manera evidente la noción de morfismo de diagramas exactos de
categorías fibradas sobre una categoría~$E$.
\end{definitions}

\subsubsection{}
\label{XIII.1.10.4}
Utilizaremos
\marginpar{359}%
en particular la noción de diagrama exacto cuando~$E$ es un sitio y
$\Phi$, $\Phi_i$, $1 \leqslant i \leqslant 3$, son stacks sobre~$E$.

Sean $f \colon E \to E'$ un morfismo de sitios y
\ifthenelse{\boolean{orig}}
{\begin{equation*}
\label{eq:XIII.1.10.4.1}
\tag{\thesubsubsection.1} { \xymatrix{ \Phi \ar[r]^p & \Phi_1
\ar@<0.6ex>[r]^{p_1} \ar@<-0.6ex>[r]_{p_2} & \Phi_2
\ar@<2ex>[r]^{p_{23}} \ar[r]^{p_{31}} \ar@<-0.6ex>[r]_{p_{12}} &
\Phi_3 } }
\end{equation*}}
{\begin{equation*}
\label{eq:XIII.1.10.4.1}
\tag{\thesubsubsection.1} { \xymatrix@C=.5cm{ \Phi \ar[r]^p & \Phi_1
\ar@<0.6ex>[rr]^{p_1} \ar@<-0.6ex>[rr]_{p_2} && \Phi_2
\ar@<2ex>[rrr]^{p_{23}} \ar[rrr]^{p_{31}} \ar@<-0.6ex>[rrr]_{p_{12}} &&&
\Phi_3 } }
\end{equation*}}
un diagrama exacto de stacks sobre~$E$. Por imagen directa se obtiene un
diagrama
$$
\xymatrix@C=.5cm{ f_* \Phi \ar[r] & f_* \Phi_1 \ar@<2pt>[r]
\ar@<-2pt>[r] & f_* \Phi_2 \ar@<4pt>[r]\ar[r] \ar@<-4pt>[r] &
f_* \Phi_3 }
$$
que evidentemente es exacto.

Si $h \colon E'' \to E$ es un morfismo de sitios, se tiene asimismo un
diagrama exacto
\ifthenelse{\boolean{orig}}
{$$
\xymatrix{ h^* \Phi \ar[r]^{p''} & h^* \Phi_1
\ar@<0.6ex>[r]^{p''_1} \ar@<-0.6ex>[r]_{p''_2} & h^* \Phi_2
\ar@<2ex>[r]^{p''_{23}} \ar[r]^{p''_{31}} \ar@<-0.6ex>[r]_{p''_{12}} &
h^* \Phi_3 }
$$}
{$$
\xymatrix@C=.5cm{ h^* \Phi \ar[r]^{p''} & h^* \Phi_1
\ar@<0.6ex>[rr]^{p''_1} \ar@<-0.6ex>[rr]_{p''_2} && h^* \Phi_2
\ar@<2ex>[rrr]^{p''_{23}} \ar[rrr]^{p''_{31}} \ar@<-0.6ex>[rrr]_{p''_{12}} &&&
h^* \Phi_3 }
$$}

Verifiquemos primero la condición~a) de~\Ref{XIII.1.10.2}. Sean
$F''$ un objeto de~$E''$, $x''$ e~$y''$ dos objetos de~$(h^*
\Phi)_{F''}$, $x''_1$ e~$y''_1$ sus imágenes respectivas en
$h^* \Phi_1$, y $x''_2$, $y''_2$ sus imágenes en~$h^*
\Phi_2$. Sea $u''_1 \colon x''_1 \to y''_1$ un morfismo tal que
$p''_1(u''_1) = p''_2(u''_1)$, y demostremos que~$u''_1$ proviene de un
único morfismo $u'' \colon x'' \to y''$. Como la cuestión es local
sobre~$F''$, podemos suponer que se tienen un objeto $F_1$ de~$E$, un
morfismo de~$F''$ en la imagen inversa $F''_1$ de~$F_1$ por~$h$, y unos
objetos $x$, $y$ de~$\Phi_{F_1}$ cuyas imágenes inversas sobre~$F''$
son $x''$ e~$y''$. Sean $x_1$, $y_1$ (\resp $x_2$, $y_2$) las imágenes
de~$x$, $y$ en~$\Phi_1$ (\resp en~$\Phi_2$). Podemos suponer que
$u''_1$ proviene de un morfismo $u_1 \colon x_1 \to y_1$ tal que
$p_1(u_1) = p_2(u_1)$. Por la exactitud de
\eqref{eq:XIII.1.10.4.1}, se obtiene un único morfismo $u \colon x
\to y$ cuya imagen inversa por~$h$ es el morfismo buscado~$u''$.

La
\marginpar{360}
condición~b) de~\Ref{XIII.1.10.2} se verifica de manera análoga. Sean
$x''_1$ un objeto de~$(h^* \Phi_1)_{F''}$ y $u''
\colon p''_1(x''_1) \to p''_2(x''_1)$ un morfismo que satisface la
relación
$$
\ifthenelse{\boolean{orig}}
{p''_{23}(u'')p''_{31}(u'')}
{p''_{31}(u'')p''_{23}(u'')}=p''_{12}(u'')^{-1} \text{,}
$$
y demostremos que existe un objeto $x''$ de~$(h^* \Phi)_{F''}$ y
un isomorfismo $i'' \colon p''(x'') \simeq x''_1$ que hace conmutativo
un diagrama análogo a~\eqref{eq:XIII.1.10.2.2}. Como la cuestión es
local sobre~$F''$, podemos suponer que se tienen un objeto $F_1$, un
morfismo $F'' \to F''_1$ como antes, y un objeto $x_1$
de~$(\Phi_1)_{F_1}$ cuya imagen inversa en~$(h^* \Phi_1)_{F''}$
es~$x''_1$. De igual manera, podemos suponer que~$u''$ proviene de un
morfismo $u \colon p_1(x_1) \to p_2(x_1)$ que satisface
\ifthenelse{\boolean{orig}}
{\eqref{eq:XIII.1.10.2.2}}
{\eqref{eq:XIII.1.10.2.1}}.
La existencia de un objeto $x$ de~$\Phi_{F_1}$ cuya imagen inversa por~$h$
sea un objeto~$x''$ que satisfaga lo requerido resulta entonces de la
exactitud de~\eqref{eq:XIII.1.10.4.1}.
\endgroup
